\chapter[Algorytmy i obliczalność]{Algorytmy, obliczalność i nierozstrzygalność}
\label{ch:computability-intro}

\section{Pytanie inne niż prawdziwość}
We wcześniejszych rozdziałach poznaliśmy zdania prawdziwe
w strukturach i zdania posiadające formalne dowody.
Teraz zadamy inne pytanie: czy istnieje \emph{jedna skończona procedura},
która dla dowolnych danych rozstrzyga określoną własność?
Możemy na przykład sprawdzić, czy dany napis jest poprawnie zbudowanym
termem języka pierwszego rzędu, analizując go od lewej do prawej.
To jest zadanie algorytmiczne. Tymczasem nie każde pytanie
o obiektach matematycznych posiada algorytm rozstrzygający.

Nie chcemy opierać się na niejasnym określeniu
„zadanie da się wykonać mechanicznie”. Wprowadzimy więc prosty model
maszyny wykonującej pojedyncze, jednoznacznie określone kroki.
Pokażemy następnie za pomocą argumentu przekątniowego,
że istnieje konkretne pytanie, na które nie odpowie
żaden algorytm działający poprawnie dla wszystkich danych.

\section{Napisy, kodowanie i problemy decyzyjne}
\begin{definicja}[Skończony alfabet i zbiór słów]
Niech $\Sigma$ będzie skończonym, niepustym zbiorem symboli.
Przez $\Sigma^*$ oznaczamy zbiór wszystkich słów skończonych
nad $\Sigma$, łącznie ze słowem pustym $\varepsilon$.
Długość słowa $w$ oznaczamy przez $|w|$.
\emph{Problemem decyzyjnym} nazywamy podzbiór $L\subseteq\Sigma^*$.
Pytanie brzmi: dla podanego słowa $w$ ustalić, czy $w\in L$.
\end{definicja}

\begin{przyklad}[Prosty problem decyzyjny]
Niech $\Sigma=\{0,1\}$, a $L$ składa się ze słów
zawierających parzystą liczbę jedynek.
Słowo $10101$ nie należy do $L$, bo ma trzy jedynki,
a słowo $1010$ należy, bo ma ich dwie.
Możemy rozstrzygnąć przynależność do $L$,
przechodząc po słowie raz i zapamiętując tylko,
czy dotychczas napotkaliśmy parzystą, czy nieparzystą liczbę jedynek.
\end{przyklad}

\begin{definicja}[Kodowanie par słów]
Dla słów binarnych $u,v$ określmy
\[
\operatorname{code}(u,v):=0^{|u|}1uv.
\]
Zapis $0^{|u|}$ oznacza ciąg $|u|$ zer stojących obok siebie;
nie jest potęgą liczby zero. Kod zaczyna się od tylu zer,
ile bitów ma pierwszy składnik, po czym zapisujemy jedynkę,
a następnie oba słowa.
\end{definicja}

\begin{przyklad}[Jak odczytać kod pary]
Dla $u=10$ i $v=01$ otrzymujemy
$\operatorname{code}(u,v)=0011001$.
Dwie początkowe zera oznaczają, że pierwszy składnik
ma długość dwa. Po pierwszej jedynce oddzielającej
odczytujemy więc dwa następne bity $10$; pozostałe bity $01$
tworzą drugi składnik. Kod daje się jednoznacznie odczytać.
\end{przyklad}

\begin{lemat}[Jednoznaczność kodowania par]\label{lem:computability-pair-code}
Dla dowolnych słów $u,v,u',v'$ nad $\{0,1\}$ zachodzi
\[
\operatorname{code}(u,v)=\operatorname{code}(u',v')
\quad\Longrightarrow\quad u=u'\ \text{oraz}\ v=v'.
\]
\end{lemat}
\begin{proof}
Każdy kod rozpoczyna się jednoznacznie określonym
blokiem zer zakończonym pierwszą jedynką.
Liczba tych zer jest długością pierwszego składnika.
Jeżeli dwa kody są równe, mają taki sam początkowy blok,
a więc $|u|=|u'|$. Kolejne $|u|$ bitów muszą być takie same,
stąd $u=u'$. Pozostała część kodu to odpowiednio $v$ i $v'$, 
więc również te słowa są równe. Opis ten daje zarazem algorytm
odkodowania: liczymy początkowe zera do pierwszej jedynki,
po czym odczytujemy tyle bitów pierwszego składnika.
Brak jedynki rozdzielającej albo zbyt mało dalszych bitów
oznacza, że napis nie jest kodem żadnej pary.
\end{proof}

\section{Maszyna Turinga: elementarne instrukcje}
\begin{definicja}[Deterministyczna maszyna Turinga]\label{def:computability-tm}
\emph{Deterministyczną maszyną Turinga} nazywamy skończony opis
złożony z:
\begin{enumerate}[label=(\roman*)]
\item skończonego zbioru stanów $Q$, zawierającego
stan początkowy $q_0$ oraz dwa różne stany końcowe
$q_{\mathrm{tak}}$ i $q_{\mathrm{nie}}$, przy czym
$q_0\notin\{q_{\mathrm{tak}},q_{\mathrm{nie}}\}$;
\item skończonego alfabetu taśmy $\Omega$ z wyróżnionym
symbolem pustym $\square$ oraz niepustym podzbiorem symboli wejściowych
$\Sigma\subseteq\Omega\setminus\{\square\}$;
\item funkcji przejścia
\[
\delta:(Q\setminus\{q_{\mathrm{tak}},q_{\mathrm{nie}}\})\times\Omega
\longrightarrow Q\times\Omega\times\{L,P,S\},
\]
której wartość $\delta(q,a)=(q',b,d)$ ma następujący sens:
„gdy maszyna jest w stanie $q$ i czyta $a$, zapisuje $b$,
przechodzi do stanu $q'$ i przesuwa głowicę
w lewo ($L$), prawo ($P$) albo pozostawia ją w miejscu ($S$)”.
\end{enumerate}
Maszyna ma taśmę o komórkach ponumerowanych liczbami
całkowitymi. Na początku wejściowe słowo jest zapisane
w kolejnych komórkach zaczynających się od komórki zero,
a wszystkie pozostałe komórki zawierają $\square$.
Głowica startuje nad komórką zero. Odczytanie stanu
$q_{\mathrm{tak}}$ oznacza \emph{akceptację}, a
$q_{\mathrm{nie}}$ \emph{odrzucenie}. W obu przypadkach
maszyna zatrzymuje się. Jeżeli nie osiągnie żadnego
ze stanów końcowych, wykonuje nieskończenie wiele kroków.
\end{definicja}

\begin{uwaga}[Ważne rozróżnienie: skończona maszyna i nieograniczony czas]
Opis maszyny, a więc wykaz stanów oraz tabela przejść,
jest skończony. Nie znaczy to jednak, że dla każdego wejścia
liczba wykonanych kroków ma ustaloną wspólną granicę.
Maszyna może poruszać się po dowolnie długich fragmentach
taśmy, a dla niektórych wejść może w ogóle się nie zatrzymać.
To właśnie możliwość nieograniczonego obliczania będzie istotna
w problemie stopu.
\end{uwaga}

\begin{przyklad}[Sprawdzenie pierwszego symbolu]
Niech alfabet wejściowy będzie równy $\{0,1\}$.
Maszyna może w stanie $q_0$ odczytać pierwszy symbol
wejściowego słowa. Jeśli jest to $1$, przechodzi
bezpośrednio do $q_{\mathrm{tak}}$, a jeśli jest to $0$
lub symbol pusty, przechodzi do $q_{\mathrm{nie}}$.
Każde obliczenie kończy się po jednym kroku.
Jest to prosty algorytm rozstrzygający własność
„słowo zaczyna się od jedynki”.
\end{przyklad}

\section{Decydowanie i rozpoznawanie}
\begin{definicja}[Zbiór rozstrzygalny i rozpoznawalny]\label{def:computability-decidable-c.e.}
Niech $L\subseteq\Sigma^*$ będzie problemem decyzyjnym.
\begin{enumerate}[label=(\roman*)]
\item Zbiór $L$ nazywamy \emph{rozstrzygalnym}, jeśli istnieje
maszyna Turinga, która zatrzymuje się dla \emph{każdego} wejścia $w$
i akceptuje dokładnie wtedy, gdy $w\in L$.
\item Zbiór $L$ nazywamy \emph{rozpoznawalnym}
(lub \emph{rekurencyjnie przeliczalnym}), jeśli istnieje
maszyna Turinga akceptująca dokładnie słowa z $L$.
Dla słów spoza $L$ maszyna może je odrzucić albo
nie zatrzymać się nigdy.
\end{enumerate}
\end{definicja}

\begin{przyklad}[Różnica między dwiema obietnicami]
Jeżeli algorytm \emph{rozstrzyga} własność,
możemy uruchomić go na dowolnych danych i otrzymać
po skończonym czasie odpowiedź „tak” lub „nie”.
Jeżeli algorytm jedynie \emph{rozpoznaje} własność,
to dla przypadku pozytywnego w końcu zobaczymy „tak”,
ale dla negatywnego możemy nie otrzymać żadnej odpowiedzi.
Z samego długiego oczekiwania nie wynika,
że odpowiedź powinna być „nie”.
\end{przyklad}

\begin{lemat}[Zbiór rozstrzygalny jest rozpoznawalny]\label{lem:computability-decidable-recognizable}
Każdy rozstrzygalny zbiór słów jest rozpoznawalny.
\end{lemat}
\begin{proof}
Niech maszyna $M$ rozstrzyga $L$.
Dla słowa należącego do $L$ maszyna zatrzymuje się
i akceptuje. Dla słowa nienależącego zatrzymuje się
i odrzuca. Są to dokładnie warunki wymagane
od maszyny rozpoznającej $L$, która może, lecz nie musi,
zatrzymywać się także na przypadkach negatywnych.
Dlatego ta sama maszyna $M$ rozpoznaje $L$.
\end{proof}

\section{Program też można zapisać jako słowo}
Każda maszyna Turinga ma skończenie wiele stanów,
symboli i wpisów w tabeli przejść. Możemy więc
zapisać jej cały opis jako skończone słowo binarne.
Ustalamy numerację stanów i symboli, zapisujemy ich liczby
oraz kolejne wpisy tabeli przejść, oddzielając pola kodami
ich długości (na przykład przez prefiks długościowy,
jak w Lemacie~\ref{lem:computability-pair-code}).
Zapis obejmuje także stan początkowy, dwa stany końcowe,
symbol pusty i alfabet wejściowy. Kodowanie to jest
jednoznacznie odczytywalne; skończona procedura sprawdza,
czy liczby symboli i stanów są poprawne, czy tabela jest
zupełna i jednoznaczna dla par $(q,a)$ ze stanu niekońcowego
oraz czy przejścia mają dozwolone kierunki.
W ten sposób każdej maszynie przypisujemy pewien napis,
który nazywamy \emph{kodem maszyny}. Przyjmujemy odtąd
jedną ustaloną, efektywną metodę takiego kodowania.
Nie każdy napis musi kodować poprawną maszynę;
niepoprawny opis rozpoznajemy podczas sprawdzania kodu.

W pozostałych rozważaniach o kodach maszyn ograniczamy się do
maszyn o alfabecie wejściowym $\{0,1\}$; wszystkie programy
konstruowane poniżej korzystają wyłącznie z takich wejść.
To wystarcza do dowodu nierozstrzygalności problemu stopu.
Zasadniczo można również kodować słowa nad innymi skończonymi
alfabetami przez słowa binarne, ustalając uprzednio ich
efektywnie odczytywalne kodowanie.

\begin{twierdzenie}[Maszyna uniwersalna]\label{thm:computability-universal}
Istnieje maszyna Turinga $U$ o binarnym alfabecie wejściowym
taka, że dla każdego poprawnego kodu $e$ maszyny $M_e$
o binarnym alfabecie wejściowym i każdego $w\in\{0,1\}^*$
maszyna $U$ uruchomiona na kodzie pary $(e,w)$
symuluje obliczenie $M_e$ na $w$ i zatrzymuje się
z taką samą odpowiedzią wtedy i tylko wtedy, gdy $M_e$
zatrzyma się na $w$.
\end{twierdzenie}
\begin{proof}
Przedstawiamy konstrukcję symulatora i uzasadniamy, dlaczego
wykonuje ją \emph{jedna} maszyna niezależna od $e$.
Najpierw $U$ odkodowuje parę $(e,w)$ według
Lematu~\ref{lem:computability-pair-code} i sprawdza poprawność
kodu $e$. Jeśli kod pary lub maszyny jest błędny, $U$ może
odrzucić wejście; te przypadki nie są objęte tezą twierdzenia.
Dla poprawnego kodu zapisuje jego tabelę przejść
oraz konfigurację początkową $M_e$ na $w$.

Konfigurację przedstawiamy skończonym napisem: obejmuje
stan, numer komórki pod głowicą (z zapisem znaku dla indeksów
ujemnych) i listę par „numer odwiedzonej komórki -- zapisany
symbol”. Indeksy zapisujemy w ustalonym kanonicznym kodzie binarnym
ze znakiem, bez nieistotnych zer wiodących; symbole kodujemy
binarnie, a kolejne pola jednoznacznie rozdzielamy.
Poza tą listą na taśmie symulowanej maszyny stoją symbole
puste. W każdej rundzie $U$ odczytuje symbol spod
symulowanej głowicy (z listy albo jako symbol pusty),
wyszukuje w zakodowanej tabeli odpowiedni wpis
$\delta(q,a)=(q',b,d)$, zmienia zapis wskazanej komórki,
stan i pozycję głowicy. Jeśli głowica wejdzie na nową
komórkę, dodaje do listy jej numer i dotychczasowy symbol pusty.

Każdą z tych czynności wykonuje się przez skończone
przeszukiwanie i przepisywanie zakodowanych napisów.
Są to operacje dostępne dla maszyny Turinga: skanowanie
skończonej listy, porównywanie zakodowanych numerów,
aktualizacja pola i obliczenie numeru sąsiedniej komórki
przez dodanie lub odjęcie jedności. W szczególności
porównanie dwóch skończonych zapisów binarnych wykonujemy
przez porównanie znaków i długości, a następnie odpowiednich
bitów, zaś dodanie lub odjęcie jedności -- przez skończony
przebieg z przeniesieniem albo pożyczką. Usunięcie lub dopisanie
pola do listy wymaga tylko przesunięcia skończonego fragmentu
zapisu na taśmie symulatora. Wszystkie te czynności można
złożyć w skończony zestaw podprogramów, niezależny od
liczby symboli zapisanych w aktualnej konfiguracji.
Kod $e$ jest \emph{danymi},
a procedura realizująca te operacje ma stałą, skończoną
tabelę instrukcji; nie zależy od konkretnej symulowanej maszyny.

Indukcyjnie po liczbie rund zapis $U$ reprezentuje dokładnie
konfigurację $M_e$ po tej samej liczbie kroków: początek
jest poprawny, a opisany odczyt i aktualizacja realizują
jedno przejście $\delta$. Jeśli $M_e$ wejdzie do stanu
$q_{\mathrm{tak}}$ lub $q_{\mathrm{nie}}$, $U$ kończy pracę
z tą samą odpowiedzią. Jeśli $M_e$ nie zatrzyma się, każda
jej kolejna konfiguracja zostanie zasymulowana,
więc także $U$ się nie zatrzyma.
\end{proof}

\begin{uwaga}[Co zostało tutaj pokazane?]
Nie twierdzimy, że każdy program ma krótki kod ani że
maszyna uniwersalna działa szybko. Istotna jest jedna
własność logiczna: istnieje jeden skończony program
potrafiący na podstawie zapisu innego programu naśladować
jego zachowanie. W szczególności programy mogą przyjmować
jako dane \emph{własne kody}.
\end{uwaga}

\section{Problem stopu}\label{sec:computability-halting}
\begin{definicja}[Problem stopu]\label{def:computability-halting}
W ustalonym binarnym kodowaniu definiujemy
$H\subseteq\{0,1\}^*$ jako zbiór słów
$\operatorname{code}(e,w)$, dla których $e$ jest poprawnym
kodem maszyny $M_e$ i maszyna $M_e$ zatrzymuje się
na wejściu $w$, niezależnie od tego, czy akceptuje,
czy odrzuca. Napisy, które nie kodują pary lub mają
niepoprawny kod maszyny, nie należą do $H$.
Pytanie problemu stopu brzmi: mając zakodowaną parę
$(e,w)$, czy możemy zawsze w skończonym czasie ustalić,
czy $M_e$ zatrzyma się na $w$?
\end{definicja}

\begin{twierdzenie}[Nierozstrzygalność problemu stopu]\label{thm:computability-halting-undecidable}
Nie istnieje maszyna Turinga rozstrzygająca zbiór $H$.
\end{twierdzenie}
\begin{proof}
Załóżmy przeciwnie, że istnieje maszyna $D$ rozstrzygająca
$H$: na każdym słowie binarnym zatrzymuje się i poprawnie
odpowiada, czy słowo to należy do $H$.
W szczególności na wejściu $\operatorname{code}(e,w)$,
gdzie $e$ jest poprawnym kodem maszyny, rozstrzyga,
czy $M_e$ zatrzyma się na $w$.

Zbudujmy program $K$ przyjmujący jako dane pojedynczy
napis $x$. Jeśli $x$ nie jest poprawnym kodem maszyny,
program $K$ może od razu zakończyć pracę.
Jeśli $x$ jest poprawnym kodem, program $K$ uruchamia
$D$ na $\operatorname{code}(x,x)$, aby ustalić,
czy maszyna o kodzie $x$ zatrzymuje się na wejściu $x$.
Następnie \emph{postępuje odwrotnie} do otrzymanej odpowiedzi:
\begin{enumerate}[label=(\roman*)]
\item jeśli $D$ przewiduje, że $M_x$ się zatrzyma,
      program $K$ zaczyna nieskończoną pętlę;
\item jeśli $D$ przewiduje, że $M_x$ się nie zatrzyma,
      program $K$ natychmiast kończy pracę.
\end{enumerate}
Każdy z tych kroków da się opisać przez maszynę Turinga:
wykonujemy obliczenie $D$ (które zawsze się kończy),
a potem albo zatrzymujemy się, albo powtarzamy jeden krok bez końca.
Zatem $K$ jest poprawnie opisanym programem i ma własny kod $k$.

Uruchommy teraz $K$ na jego własnym kodzie $k$.
Jeśli $D$ na wejściu $\operatorname{code}(k,k)$ odpowie,
że $M_k=K$ zatrzyma się na $k$, to z definicji
$K$ zapętli się i \emph{nie} zatrzyma.
Jeśli $D$ odpowie, że $K$ \emph{nie} zatrzyma się na $k$,
to $K$ natychmiast się zatrzyma.
W obu przypadkach przewidywanie $D$ jest błędne.
To przeczy założeniu, że $D$ zawsze udziela poprawnej odpowiedzi.
Zatem taki program nie istnieje.
\end{proof}

\begin{figure}[htbp]
\centering
\begin{tikzpicture}[
  >=Stealth,
  font=\small,
  pytanie/.style={draw,rounded corners,align=center,minimum width=38mm,minimum height=9mm},
  przypadek/.style={draw,rounded corners,align=center,text width=40mm,minimum height=16mm},
  line width=.55pt]
\node[pytanie] (d) {$D(\operatorname{code}(k,k))$};
\node[przypadek,below left=11mm and 8mm of d] (tak)
  {$D$ mówi: „zatrzyma się”\\ $K(k)$ nie zatrzymuje się};
\node[przypadek,below right=11mm and 8mm of d] (nie)
  {$D$ mówi: „nie zatrzyma się”\\ $K(k)$ zatrzymuje się};
\draw[->] (d.south west) -- (tak.north east);
\draw[->] (d.south east) -- (nie.north west);
\end{tikzpicture}
\caption{Argument przekątniowy: w obu możliwych przypadkach
  odpowiedź $D$ jest sprzeczna z zachowaniem programu $K$.}
\label{fig:computability-diagonal}
\end{figure}

\begin{uwaga}[Dlaczego to nie jest paradoks językowy?]
W dowodzie nie zakładaliśmy, że jakiś program
ma jednocześnie się zatrzymać i nie zatrzymać.
Przypuściliśmy jedynie istnienie \emph{uniwersalnego
rozstrzygającego} programu $D$. Na tej podstawie
skonstruowaliśmy nowy program $K$, którego zachowanie
przeczyłoby odpowiedzi $D$ dla jednego wybranego wejścia.
Wniosek jest zatem taki, że \emph{nie istnieje} program
$D$ działający poprawnie we wszystkich przypadkach.
\end{uwaga}

\section{Problem stopu jest mimo to rozpoznawalny}
\begin{twierdzenie}[Rozpoznawalność problemu stopu]\label{thm:computability-halting-recognizable}
Zbiór $H$ jest rozpoznawalny, ale nie jest rozstrzygalny.
\end{twierdzenie}
\begin{proof}
Nierozstrzygalność udowodniliśmy w
Twierdzeniu~\ref{thm:computability-halting-undecidable}.
Aby rozpoznać $H$, bierzemy dowolne słowo binarne.
Sprawdzamy, czy jest kodem pary $(e,w)$; jeśli nie,
odrzucamy. Następnie sprawdzamy, czy $e$ opisuje
poprawną maszynę; w przeciwnym razie także odrzucamy.
Dla poprawnego kodu symulujemy działanie $M_e$ na $w$
za pomocą maszyny uniwersalnej z
Twierdzenia~\ref{thm:computability-universal}.
Jeżeli symulowana maszyna zatrzyma się,
rozpoznający akceptuje. Jeżeli się nie zatrzyma,
symulacja trwa w nieskończoność, ale definicja
rozpoznawalności na to pozwala.
\end{proof}

\begin{twierdzenie}[Rozstrzygalność przez dwa rozpoznawania]\label{thm:computability-two-recognizers}
Dla dowolnego $L\subseteq\{0,1\}^*$ zbiór $L$
jest rozstrzygalny wtedy i tylko wtedy, gdy zarówno
$L$, jak i jego dopełnienie $\{0,1\}^*\setminus L$
są rozpoznawalne.
\end{twierdzenie}
\begin{proof}
Jeśli $L$ jest rozstrzygalny, to rozpoznajemy $L$
maszyną rozstrzygającą. Aby rozpoznać dopełnienie,
wystarczy zamienić jej odpowiedzi „tak” i „nie”.

Odwrotnie, niech $R_1$ rozpoznaje $L$, a $R_0$
rozpoznaje jego dopełnienie. Dla zadanego słowa $w$
nie uruchamiamy $R_1$ i nie oczekujemy bez końca,
a dopiero potem $R_0$. Zamiast tego przeplatamy obliczenia:
wykonujemy jeden krok $R_1$, jeden krok $R_0$,
następny krok $R_1$, następny krok $R_0$ itd.
Jeżeli któryś z programów zaakceptuje $w$,
zatrzymujemy się i zwracamy odpowiadającą mu decyzję.
Ewentualne odrzucenie przez jedną z maszyn ignorujemy:
z punktu widzenia rozpoznawalności negatywny wynik
nie jest świadectwem przynależności do jej zbioru.
Zatrzymanej bez akceptacji symulacji nie wznawiamy,
a drugą kontynuujemy.
Każde $w$ należy albo do $L$, albo do jego dopełnienia,
więc przynajmniej jeden rozpoznający musi je po skończonej
liczbie kroków zaakceptować. Opisana procedura
zawsze się zatem zatrzymuje i poprawnie rozstrzyga $L$.
\end{proof}

\begin{wniosek}[Dopełnienie problemu stopu nie jest rozpoznawalne]
Zbiór $\{0,1\}^*\setminus H$ nie jest rozpoznawalny.
\end{wniosek}
\begin{proof}
Gdyby dopełnienie $H$ było rozpoznawalne,
to na mocy Twierdzenia~\ref{thm:computability-two-recognizers}
zbiór $H$ byłby rozstrzygalny, ponieważ sam $H$
jest rozpoznawalny. Przeczy to
Twierdzeniu~\ref{thm:computability-halting-undecidable}.
\end{proof}

\section{Dowody formalne jako dane dla algorytmu}
Zwróćmy uwagę na ważne powiązanie z rozdziałem~\ref{ch:natural-deduction}.
Poprawność pojedynczego kroku dedukcji naturalnej
jest sprawdzalna, ponieważ mamy skończony zapis reguły
oraz jej warunków ubocznych. Przy ustalonym efektywnie
opisanym, przeliczalnym języku i obliczalnym systemie aksjomatów
możemy więc przeglądać kolejno skończone drzewa
i sprawdzać, które są dowodami.

\begin{twierdzenie}[Rozpoznawalność dowodliwości w efektywnym systemie]
\label{thm:computability-proof-recognizable}
Niech $\mathcal L$ będzie językiem pierwszego rzędu,
w którym poprawność kodów symboli, ich rodzaje i arności
można rozstrzygać algorytmicznie. Zakładamy, że termy,
formuły oraz zdania języka $\mathcal L$ mają skończone
kody binarne, których poprawność można sprawdzać
algorytmicznie. Niech $T$ będzie rozstrzygalnym zbiorem zdań
tego języka, a sprawdzanie wszystkich reguł dedukcji naturalnej
(wraz z warunkami ubocznymi) niech będzie algorytmiczne.
Wtedy zbiór kodów zdań $\sigma$ spełniających
$T\vdash_{\mathrm{ND}}\sigma$ jest rozpoznawalny.
\end{twierdzenie}
\begin{proof}
Na wejściu najpierw sprawdzamy, czy otrzymany napis
jest kodem zdania $\sigma$; w przeciwnym razie go odrzucamy.
Następnie wypisujemy kolejno wszystkie skończone słowa
binarne, na przykład według długości i leksykograficznie.
Każdy dowód formalny jest skończonym drzewem o skończonym
kodzie binarnym, więc jego kod pojawi się w tym przeglądzie.
Dla każdego słowa sprawdzamy,
czy koduje poprawne drzewo: czy wszystkie węzły
zawierają poprawne formuły, czy każdy krok używa
jednej z dozwolonych reguł z zachowaniem warunków ubocznych
i czy wszystkie pozostałe otwarte założenia należą do $T$.
Ostatni warunek jest rozstrzygalny z założenia o $T$;
pozostałe kontrole wynikają z jawnych założeń o składni
i regułach.
Jeżeli znajdziemy drzewo z wnioskiem $\sigma$, akceptujemy.
Jeśli takie drzewo istnieje, będzie jednym ze skończonych
napisów w naszym przeglądzie, więc w końcu je znajdziemy.
Jeśli nie istnieje, przegląd może trwać bez końca;
to dokładnie odpowiada definicji rozpoznawalności.
\end{proof}

\begin{przyklad}[Nie mylmy sprawdzenia dowodu ze znalezieniem dowodu]
Gdy ktoś poda nam \emph{konkretne} drzewo dowodu,
możemy krok po kroku skontrolować jego poprawność.
Nie oznacza to jednak, że istnieje algorytm,
który dla dowolnego zdania zawsze rozstrzygnie,
czy \emph{jakieś} poprawne drzewo dowodu istnieje.
Sprawdzanie gotowego certyfikatu i rozstrzyganie
istnienia certyfikatu to dwa różne problemy.
\end{przyklad}


\section{Przejście od maszyn Turinga do rachunku na liczbach}
W rozdziale~\ref{ch:arithmetization} wprowadzimy funkcje
pierwotnie rekurencyjne. Są one szczególną klasą funkcji
obliczalnych, ale nie należy utożsamiać tej klasy ze wszystkimi
funkcjami obliczalnymi przez maszyny Turinga.
W tym miejscu chodzi o mechanizm kodowania: liczby naturalne
mogą reprezentować słowa i konfiguracje maszyny, a pojedynczy
krok obliczenia daje się opisać działaniami arytmetycznymi.

\begin{przyklad}[Dwa fragmenty taśmy jako dwie liczby]
Niech alfabet taśmy ma $B\geqslant 2$ symboli oznaczonych
liczbami $0,\ldots,B-1$, przy czym symbolowi pustemu
przypisujemy numer $0$. Niepusty skończony napis
$a_0a_1\cdots a_k$ przedstawiamy liczbą
$A=a_0+a_1B+\cdots+a_kB^k$ i osobno zapisujemy jego
długość $k+1$. Pusty napis reprezentujemy przez parę $(0,0)$.
Dla napisu niepustego pierwszy symbol jest równy $A\bmod B$,
a kod pozostałej części wynosi $\lfloor A/B\rfloor$.

Zakodujmy tak dwa skończone fragmenty taśmy, leżące po lewej
i po prawej stronie głowicy, w kolejności \emph{od komórek
najbliższych głowicy}. Niech ich kody i długości będą równe
$(L,\ell)$ i $(R,r)$, a $a$ oznacza symbol aktualnie czytany.
Załóżmy, że funkcja przejścia nakazuje
$\delta(q,a)=(q',c,d)$, a więc maszyna \emph{zapisuje symbol $c$}
(w ogólności różny od $a$) i wykonuje ruch $d$.

Jeśli $d=P$, to zapisany symbol $c$ staje się pierwszym
symbolem po lewej stronie nowego położenia głowicy.
Otrzymujemy $L'=c+BL$ i $\ell'=\ell+1$.
Jeśli $r>0$, nowo odczytany symbol jest równy
$a'=R\bmod B$, zaś $R'=\lfloor R/B\rfloor$ i $r'=r-1$.
Jeśli $r=0$, otrzymujemy $a'=0$, $R'=0$ i $r'=0$.

Analogicznie, jeśli $d=L$, to $R'=c+BR$ i $r'=r+1$.
Dla $\ell>0$ otrzymujemy $a'=L\bmod B$,
$L'=\lfloor L/B\rfloor$ i $\ell'=\ell-1$,
a dla $\ell=0$ otrzymujemy $a'=0$, $L'=0$ i $\ell'=0$.
Wreszcie jeśli $d=S$, głowica się nie przesuwa:
$L'=L$, $\ell'=\ell$, $R'=R$, $r'=r$ oraz $a'=c$.
We wszystkich trzech przypadkach nowy stan wynosi $q'$.

Wzory te realizują zapis i ruch głowicy także wtedy,
gdy głowica przechodzi na dotąd nieodwiedzoną komórkę.
Oddzielne długości zapobiegają utracie końcowych zer
w kodach fragmentów taśmy. W ten sposób \emph{cały krok}
maszyny (zapis, przesunięcie i zmiana stanu) zostaje wyrażony
przez elementarne operacje arytmetyczne.
\end{przyklad}

\begin{uwaga}[Znaczenie dla dalszych dowodów]
Weryfikacja \emph{podanego} skończonego śladu maszyny
różni się od wyszukiwania śladu zakończonego zatrzymaniem.
W pierwszym przypadku wystarczy skończona liczba
kontroli sąsiednich konfiguracji. W drugim liczba
prób nie ma z góry ograniczenia. Właśnie ta różnica
umożliwi w dalszej części jednoczesne kodowanie
sprawdzalności dowodu i nierozstrzygalności dowodliwości.
\end{uwaga}

\section{Związek z twierdzeniami Gödla}
Twierdzenie o pełności~\ref{thm:godel-completeness} utożsamia
konsekwencję semantyczną z dowodliwością w logice pierwszego rzędu.
Nie zapewnia jednak algorytmu, który dla każdej formuły
rozstrzyga istnienie dowodu. Systematyczny przegląd dowodów
wykrywa dowód, jeśli istnieje; w pozostałych przypadkach
może nie zakończyć pracy.

W dalszych rozdziałach zakodujemy formuły i dowody liczbami
naturalnymi oraz wykorzystamy to w twierdzeniach Gödla
o niezupełności. Przed włączeniem istniejącego opracowania
sprawdzimy jego dowody i spójność notacji.